\begin{proofbox}
\textbf{\textcolor{CProof}{思路提示}} 令 $t=f(x)$，始终假设 $M>N$ 且 $x\in D$。保留原文四步转换路线，每步分别验证正向和反向。
\begin{description}[style=nextline,leftmargin=2.8em,labelsep=0.5em,itemsep=0.55em]
\item[\textbf{步骤一（乘积形式）：}]
\textbf{正向：} 从 $N<t<M$ 出发，将其拆成 $t>N$ 与 $t<M$。
分别减去 $N$ 和 $M$，得到
\[t-N>0,\qquad t-M<0.\]
一正一负相乘为负，因此
\[(t-M)(t-N)<0.\]
\textbf{反向：} 若 $(t-M)(t-N)<0$，两因式异号，只有两种情况：
\[
\begin{cases}t-M<0,\\t-N>0\end{cases}
\quad\text{或}\quad
\begin{cases}t-M>0,\\t-N<0.\end{cases}
\]
第一种分别移项得 $t<M$、$t>N$，即 $N<t<M$。
第二种分别移项得 $t>M$、$t<N$，于是 $M<t<N$，推出 $M<N$，与已知 $M>N$ 矛盾。
故只能是第一种，得到
\[N<t<M\iff(t-M)(t-N)<0.\]
\item[\textbf{步骤二（绝对值形式）：}]
\textbf{正向：} 从乘积形式配方。先展开：
\begin{align*}
(t-M)(t-N)&=t^2-Nt-Mt+MN\\
&=t^2-(M+N)t+MN.
\end{align*}
一次项系数的一半为 $-(M+N)/2$，其平方为 $(M+N)^2/4$。同时加上并减去这个数，不改变原式：
\begin{align*}
(t-M)(t-N)&=t^2-(M+N)t\\
&\quad+\frac{(M+N)^2}{4}-\frac{(M+N)^2}{4}+MN\\
&=\left(t-\frac{M+N}{2}\right)^2
-\frac{M^2+2MN+N^2-4MN}{4}\\
&=\left(t-\frac{M+N}{2}\right)^2
-\frac{M^2-2MN+N^2}{4}\\
&=\left(t-\frac{M+N}{2}\right)^2
-\left(\frac{M-N}{2}\right)^2.
\end{align*}
令 $C=(M+N)/2$，$R=(M-N)/2$。因为 $M>N$，所以 $R>0$。
若乘积为负，将上述恒等式代入：
\[(t-C)^2-R^2<0.\]
两边加上 $R^2$，得 $(t-C)^2<R^2$。
因为 $|t-C|\ge0$ 且 $R>0$，平方在非负数范围内保持大小关系，所以开平方得到 $|t-C|<R$。

\textbf{反向：} 若 $|t-C|<R$，两边均非负，平方得
\[|t-C|^2<R^2.\]
利用 $|t-C|^2=(t-C)^2$，再两边减去 $R^2$，得到
\[(t-C)^2-R^2<0.\]
代入刚才的配方恒等式，得到 $(t-M)(t-N)<0$。故两形式等价。
\item[\textbf{步骤三（分式形式）：}]
\textbf{正向：} 由 $N<t<M$ 分别移项，得 $t-N>0$ 和 $M-t>0$。
分母严格为正，故不为零；正数除以正数仍为正：
\[\frac{t-N}{M-t}>0.\]
\textbf{反向：} 若上式成立，则 $M-t\ne0$，且分子分母同号。设 $A=t-N$，$B=M-t$，则
\[A+B=t-N+M-t=M-N>0.\]
若 $A<0,B<0$，则 $A+B<0$，与上式矛盾。因此两者只能同时为正，即
\[t-N>0,\qquad M-t>0.\]
分别移项得 $t>N$、$t<M$，合并即 $N<t<M$。故两形式等价。
\item[\textbf{步骤四（倒数形式）：}]
\textbf{正向（保留分式过渡）：} 由分式形式及步骤三，知 $t-N>0$、$M-t>0$。
将正分式取倒数，仍为正：
\[\frac{M-t}{t-N}>0.\]
分子写成 $(M-N)-(t-N)$，于是
\begin{align*}
\frac{(M-N)-(t-N)}{t-N}&>0\\
\frac{M-N}{t-N}-1&>0\\
\frac{M-N}{t-N}&>1\\
\frac1{t-N}&>\frac1{M-N}.
\end{align*}
最后一步除以正数 $M-N$，不等号方向不变。

\textbf{正向（原文直接方法）：} 由 $N<t<M$，三部分同时减去 $N$：
\[0<t-N<M-N.\]
两个量 $t-N$ 与 $M-N$ 都为正数。函数 $y=1/u$ 在 $u>0$ 时递减，故取倒数后大小关系反转：
\[\frac1{t-N}>\frac1{M-N}.\]
\textbf{反向：} 若倒数不等式成立，由 $M-N>0$ 知右边为正，故左边也为正。这说明 $t-N>0$（且倒数存在）。两个正数取倒数后大小关系反转，所以
\[t-N<M-N.\]
合并得 $0<t-N<M-N$，三部分同时加上 $N$，得到 $N<t<M$。
因此原文的正值条件可以保留作为操作提醒，但它并非额外假设。

\textbf{补充证明（作差通分）：} 由 $N<t<M$ 知 $t-N>0$、$M-t>0$，且 $M-N>0$。
\begin{align*}
\frac1{t-N}-\frac1{M-N}
&=\frac{(M-N)-(t-N)}{(t-N)(M-N)}\\
&=\frac{M-t}{(t-N)(M-N)}>0.
\end{align*}
分子、分母均为正，故差为正，即所证。
反之，倒数不等式说明 $1/(t-N)>1/(M-N)>0$，所以 $t-N>0$。通分后上式的分母为正；由于差为正，分子 $M-t>0$。于是 $N<t<M$。
\end{description}
\medskip
\textbf{\textcolor{CProof}{结论回扣}}
\begin{center}
\begin{tikzpicture}[every node/.style={draw,rounded corners,align=center,font=\small}]
\node (a) at (0,0){原始形式};
\node (b) at (-2.3,-1.3){乘积形式};
\node (c) at (-2.3,-2.6){绝对值形式};
\node (d) at (2.3,-1.3){分式形式};
\node (e) at (2.3,-2.6){倒数形式};
\draw[<->] (a)--(b);\draw[<->] (b)--(c);\draw[<->] (a)--(d);\draw[<->] (d)--(e);
\end{tikzpicture}
\end{center}
在 $M>N$ 且函数值为实数的条件下，五种形式彼此等价；每条转换均可逆。
\end{proofbox}
